\section{Sampling algorithms and exact correction}\label{sec:algorithms}

We now turn the weighted graph bound into an algorithm using unbiased random
bits. All deterministic choices, including greedy completions, neighbor
orders, and tie breaking, are fixed once and for all. Abbreviate
$\Omega(r,c)$ by $\Omega$. Write
$\nu_X$ for the uniform law on $\mathcal F_X$ and $z_0$ for a greedy original
table. Its restriction to the small slots gives a starting transversal $X_0$.

\subsection{The ideal chain}

Set
\[
 q_{\rm s}=\lceil\log_2(32d^2)\rceil,
 \qquad p=2^{-q_{\rm s}}.
\]
List the distinct feasible neighbors of a state $X$ in a fixed order.
Assign one of the $2^{q_{\rm s}}$ equally likely binary strings to each
neighbor and let all remaining strings propose a hold. This gives every
edge the same proposal probability $p$ in both directions. As explained in
Lemma~\ref{lem:repair}, the list is obtained by testing all unit exchanges,
the repair of $X$ when it is a defect, and the uniquely reconstructed inverse
repair of each possible ordered defect type. There are at most $5d^2+1$
neighbors. Each test checks only coordinate bounds, the allowed occupancy
pattern, and residual margins; it does not evaluate $f$.

Here is the completion adjustment used for a nonholding proposal $X\to Y$.
For a nonempty block $\mathcal B=I\times J$, fix a reference row $i_0$ and
column $j_0$, independently of $X$. Let
$\Delta R_i=R_i(Y)-R_i(X)$ and $\Delta P_j=P_j(Y)-P_j(X)$ be the changes in
its padded margins. Define the integer matrix $D_{XY}$ by
\[
 (D_{XY})_{ij}=\begin{cases}
 \Delta R_i,&i\ne i_0,\ j=j_0,\\
 \Delta P_j,&i=i_0,\ j\ne j_0,\\
 \Delta R_{i_0}-\sum_{j'\ne j_0}\Delta P_{j'},&i=i_0,\ j=j_0,\\
 0,&i\ne i_0,\ j\ne j_0.
 \end{cases}
\]
The equality of total row and column changes verifies all its row and column
sums, including the reference column. Also $D_{YX}=-D_{XY}$. Along a unit
exchange or repair edge each individual residual changes by at most two;
hence every displayed entry has magnitude at most $10d$. These formulas
also apply to a block with a single row or column.

The ideal transition draws $a\sim\nu_X$ and accepts $Y$ precisely when
$a+D_{XY}$ is nonnegative. It then belongs to $\mathcal F_Y$ because its
margins are already correct. A rejected proposal holds at $X$. If the block
is empty, every edge proposal is accepted. Denote this ideal kernel by $P$.

\begin{proposition}[Small-chain bound]\label{prop:small-chain}
The kernel $P$ is lazy and reversible for $\pi(X)=f(X)/\Lambda$. On every
edge it satisfies
\[
 P(X,Y)\ge \frac{1}{128d^2}
                  \min\left\{1,\frac{f(Y)}{f(X)}\right\}.
\]
If there is more than one state, its inverse spectral gap is at most
$d^{160}$. Moreover
$\pi_{\min}\ge(C+1)^{-2d}$, and, for every $t\ge0$,
\begin{equation}\label{eq:small-mixing}
 \TV\bigl(P^t(X_0,\cdot),\pi\bigr)
 \le (C+1)^{2d}\exp(-t/d^{160}).
\end{equation}
A one-state chain is already stationary.
\end{proposition}

\begin{proof}
The proposal holds with probability at least
$1-(5d^2+1)/(32d^2)>1/2$, and rejection can only increase this probability.
Translation by $D_{XY}$ bijects the set of acceptable completions in
$\mathcal F_X$ with the corresponding set in $\mathcal F_Y$. Calling its
cardinality $c_{XY}=c_{YX}$ gives
$f(X)P(X,Y)=p c_{XY}=f(Y)P(Y,X)$.

Every nonempty completion-block margin is at least $L$. Rejection implies
that one of at most $d$ entries of a uniform completion is less than $10d$.
Lemma~\ref{lem:small-entry}, with $a=L$ and threshold $10d\le L/2$, gives
\[
 \Pp(\text{rejection})\le \frac{40d^5}{L}
             =\frac{40}{d^7}<\frac12.
\]
Consequently $c_{XY}\ge f(X)/2$, and in particular
$c_{XY}\ge\min(f(X),f(Y))/2$. Since $p\ge1/(64d^2)$, both the stated
transition bound and the Dirichlet-form comparison
\[
 \mathcal E_P(H)\ge \frac{E(H)}{128d^2\Lambda}
\]
follow. Theorem~\ref{thm:poincare} and $U=d^{20}$ now give an inverse-gap
bound of
\begin{align*}
128d^2\bigl((1+2d^2)4d^2(U+1)^6+2\bigr)
 &\le 98{,}560d^{126}\le d^{160}.
\end{align*}
Indeed the coefficient inside parentheses is at most $770d^4U^6$.
The graph is connected by Theorem~\ref{thm:poincare}, and
Proposition~\ref{prop:stationary-success} gives
$\pi_{\min}\ge(C+1)^{-2d}$. Lemma~\ref{lem:mixing} now gives
\eqref{eq:small-mixing}.
\end{proof}

\subsection{A bounded almost-uniform sampler}

For an integer $k\ge1$, put
\begin{equation}\label{eq:outer-parameters}
 J_{\rm o}=d^4(k+1),\qquad
 T=d^{200}(k+b)^2,\qquad h=d^4(k+b)^2.
\end{equation}
The implemented chain replaces every ideal completion draw by the
subroutine of Theorem~\ref{thm:dense-draw} with accuracy $h$; write $\mu_X$
for its output law on $\mathcal F_X$ and $Q$ for the resulting kernel.
Starting anew at $X_0$ on each trial, perform the following bounded procedure.
\begin{enumerate}
\item Run $Q$ for $T$ transitions and then make one further, independent
completion draw with law $\mu_X$ at the final state.
\item If $X$ is a transversal and every large completion entry is at least
$L$, subtract $L$ on the large block and return the resulting original
table. Otherwise this trial fails.
\item Repeat for at most $J_{\rm o}$ independent trials, and return $z_0$
if all of them fail.
\end{enumerate}
All randomness in different draws and trials is fresh. In particular, a
completion used to test a transition is discarded after that test; the
current small state is the chain's entire persistent state.

\begin{theorem}[Almost-uniform sampling]\label{thm:almost-uniform}
The preceding procedure returns a feasible table with law $p_k$ satisfying
\[
 \TV\bigl(p_k,\operatorname{Uniform}(\Omega)\bigr)\le2^{-k}.
\]
Its worst-case number of bit operations is bounded by a fixed polynomial
in $d,b,k$.
\end{theorem}

\begin{proof}
We first control adaptive approximation. The dense guarantee is uniform
over feasible states:
$\TV(\mu_X,\nu_X)\le\eta:=2^{-h}$.
From a common state, couple the proposal and holding decision identically,
and maximally couple the two completion draws if there is a nonholding
proposal. Agreement of these draws gives the same acceptance decision and
the same next state. Thus
\[
 \sup_X\TV\bigl(Q(X,\cdot),P(X,\cdot)\bigr)\le\eta.
\]
Couple successively until the first disagreement. Conditional on agreement
so far, the common state may be adaptive, but the same uniform conditional
bound still applies. After $T$ transitions and the terminal completion, the
probability of a disagreement is at most $(T+1)\eta$. This argument uses
neither reversibility nor a stationary-law assertion for $Q$.
By Proposition~\ref{prop:small-chain}, the additional error from replacing
the ideal terminal state by a stationary one is at most
$\tau=(C+1)^{2d}e^{-T/d^{160}}$.

In a stationary ideal trial every pair $(X,a)$ has probability
$\pi(X)/f(X)=1/\Lambda$. By Proposition~\ref{prop:stationary-success},
successful pairs are in bijection with $\Omega$, and the success probability
$s_*=|\Omega|/\Lambda$ is at least $1/[2(1+d^2)]$.
The first success among $J_{\rm o}$ independent stationary trials therefore
has, including the fallback, the exact law
\[
 \bigl[1-(1-s_*)^{J_{\rm o}}\bigr]
          \operatorname{Uniform}(\Omega)
       +(1-s_*)^{J_{\rm o}}\delta_{z_0}.
\]
For clarity, each particular table receives first-success mass
$\sum_{j=0}^{J_{\rm o}-1}(1-s_*)^j/\Lambda
=[1-(1-s_*)^{J_{\rm o}}]/|\Omega|$.
Pregenerate all potential trials in both procedures, then apply the same
deterministic first-success-or-fallback map. The union bound and contraction
of total variation under a deterministic map give total error at most
\begin{equation}\label{eq:outer-error}
 e^{-J_{\rm o}/[2(1+d^2)]}
 +J_{\rm o}(C+1)^{2d}e^{-T/d^{160}}
 +J_{\rm o}(T+1)2^{-h}.
\end{equation}
In particular, no error is divided by a success probability.

Each term in \eqref{eq:outer-error} is at most $2^{-k-2}$. Here are explicit
checks. Put $t=k+b\ge15$, so $d\le t$, $k+2\le t$, and
$\log_2(C+1)<b$. First,
\[
 \frac{J_{\rm o}}{2(1+d^2)}
 \ge\frac{d^2(k+1)}4\ge49(k+1)\ge(k+2)\ln2.
\]
For the second term,
\[
 \ln\bigl(J_{\rm o}(C+1)^{2d}\bigr)+(k+2)\ln2
 \le2t^2+6t\le3t^2\le d^{40}t^2=T/d^{160}.
\]
Finally, since $T+1\le2T$,
\begin{align*}
 \log_2\bigl(J_{\rm o}(T+1)\bigr)+k+2
 &\le204\log_2d+2\log_2t+\log_2(k+1)+k+3\\
 &\le209t\le d^4t^2=h.
\end{align*}
Their sum is less than $2^{-k}$.

For the running time, neighbor lists have polynomial size, and all state
coordinates, residuals, and completion entries have $O(b+\log d)$ bits.
The proposal uses exactly $q_{\rm s}$ unbiased bits. Each completion call
has bounded polynomial bit cost by Theorem~\ref{thm:dense-draw}; the numbers
of steps, trials, and calls in \eqref{eq:outer-parameters} are fixed
polynomials. Comparisons, additions, integer multiplications and divisions,
and loop counters all have polynomial bit cost, for example with schoolbook
arithmetic. Only feasible completions and the feasible fallback can be
returned. Empty and singleton-dimensional blocks use their deterministic
completion rule. This proves the claimed worst-case bound.
\end{proof}

\subsection{Exact tabulation of the implemented law}

The exact correction will need the law of the actual bounded algorithm,
including approximate offsets, adaptive transitions, and both stopping
rules. Only the small state persists between transitions. We can therefore
cache the dense completion law for each small state, propagate the resulting
finite kernel, and then apply the terminal and first-success rules. The spaces
enumerated in this calculation do not grow with the accuracy request;
accuracy changes only the iteration counts and the lengths of integers.

\begin{proposition}[Law tabulation]\label{prop:law-tabulation}
For fixed margins and $k$, the entire law $p_k$ can be computed exactly in
\[
 2^{500db}\operatorname{poly}(d,b,k)
\]
bit operations. It has a common dyadic denominator:
$p_k(z)=a_z/2^R$, where $a_z$ are nonnegative integers and $R$ is bounded by
a fixed polynomial in $d,b,k$. The sizes of all enumerated state spaces
are independent of $k$.
\end{proposition}

\begin{proof}
\emph{The finite spaces.}
Use the dense sampler's notation $A=d^4$ and $K_0=4d^8$.
There are at most $\mathsf S:=2^{2db}$ feasible small states, since their
number is at most $(C+1)^{2d}$. A completion fiber, the original table
space, an offset space, and the bin space each have at most
$\mathsf V:=2^{db}$ elements: use $(C+1)^d\le2^{db}$ and $K_0<C$.
They can be enumerated by listing coordinate arrays and retaining those
that pass their defining tests. None of these coordinate ranges depends
on $k$. This exhaustive calculation is used only for the exact correction,
not during an ordinary run of the bounded sampler.

\emph{Caching a dense completion law.}
The following denominator allowances are uniform over all feasible small
states, with $h$ as in \eqref{eq:outer-parameters}:
\begin{align}
 Q_0&=\lceil\log_2(8d)\rceil+8Ad,
 &r_{\max}&=b+1+4(h+d),\nonumber\\
 H&=d^{50}(h+1)^2,
 &J_{\rm D}&=d^4(h+1),\nonumber\\
 R_1&=H Q_0+d r_{\max},
 &R_{\rm D}&=J_{\rm D}R_1.\label{eq:bit-schedule}
\end{align}
Every bin transition has denominator dividing $2^{Q_0}$: its proposal
needs $\lceil\log_2(8d)\rceil$ bits and every dyadic acceptance test needs
at most $8Ad$ bits. For each input $X$, enumerate its bins, compute their
transition matrix as $B_X/2^{Q_0}$ with integer entries, and propagate the
point mass at zero for $H$ steps. At step $t$ the vector uses denominator
$2^{tQ_0}$; its next numerator is the previous integer vector multiplied
by $B_X$. Adding contributions from many paths changes the numerator,
not the denominator exponent.

Offsets are tabulated without enumerating their random bits. For width
$w$ and the actual bit count
$r=\lceil\log_2(w+1)\rceil+4(h+d)\le r_{\max}$, the probability of offset
$u\in\{0,\ldots,w-1\}$ is
\begin{equation}\label{eq:offset-count}
 2^{-r}\left(
 \left\lceil\frac{(u+1)2^r}{w}\right\rceil
 -\left\lceil\frac{u2^r}{w}\right\rceil\right).
\end{equation}
Indeed these two endpoints count the integers $z$ with
$u\le wz/2^r<u+1$. Integer ceiling division evaluates the count exactly.
Enlarge its denominator to $2^{r_{\max}}$, pad unused free coordinates to
$d$, and combine with the propagated bin law. Enumerating bin and offset
tuples, constructing the candidate table, and applying its feasibility
test gives the implemented dense trial's success numerators $\alpha_X(a)$ and failure
numerator $\phi_X$, all with denominator $2^{R_1}$.
Sorting of margins and its inverse are deterministic functions of $X$;
we always record completions in the original block coordinates.

Let $a_X^0$ be the dense call's fixed greedy fallback. The complete dense
output law is exactly
\begin{equation}\label{eq:cached-dense-law}
 \mu_X(a)=\frac{
 \displaystyle\sum_{j=0}^{J_{\rm D}-1}
    \phi_X^j\alpha_X(a)2^{R_1(J_{\rm D}-j-1)}
    +\one_{\{a=a_X^0\}}\phi_X^{J_{\rm D}}
 }{2^{R_{\rm D}}}.
\end{equation}
This is a finite first-success sum and the all-fail mass. It introduces
no division by a success probability. A deterministic empty or
singleton-dimensional completion call uses the same denominator by
ignoring the reserved bits. Cache \eqref{eq:cached-dense-law} for each $X$.

\emph{Propagating the implemented chain.}
For a distinct neighbor $Y$ the implemented transition is
\begin{equation}\label{eq:implemented-kernel}
 \begin{aligned}
 Q(X,Y)=p\sum_{a\in\mathcal F_X}\mu_X(a)
                \one_{\{a+D_{XY}\in\mathcal F_Y\}},\\
 Q(X,X)=1-\sum_{Y\ne X}Q(X,Y).
 \end{aligned}
\end{equation}
For an empty block interpret the indicator as one. These entries share
denominator $2^{R_{\rm D}+q_{\rm s}}$. Since only $X$ persists between
steps and completion randomness is fresh, ordinary vector propagation
$v_{t+1}=v_tQ$, starting at $X_0$, computes the exact adaptive state law.
Append the terminal completion draw and the unpadding test. This yields
the one-outer-trial success masses $g(z)$ and failure mass
$f_0=1-\sum_z g(z)$ with common denominator $2^{R_{\rm T}}$, where
\[
 R_{\rm T}=T(R_{\rm D}+q_{\rm s})+R_{\rm D}.
\]
Explicitly, $g(z)$ is the sum of $v_T(X)\mu_X(a)$ over pairs passing that
test and producing $z$. The final law is
\begin{equation}\label{eq:outer-geometric-law}
 p_k(z)=g(z)\sum_{j=0}^{J_{\rm o}-1}f_0^j
                 +\one_{\{z=z_0\}}f_0^{J_{\rm o}},
\end{equation}
whose denominator divides $2^R$ for $R=J_{\rm o}R_{\rm T}$.
Both geometric sums are evaluated by bounded integer multiplication and
addition. Algebraically, the common denominators reserve ignored bits for
holds, shorter draws, and stopped trials; neither tabulation nor sampling
enumerates those bit strings.

\emph{Bit complexity.}
The denominator exponents are bounded by fixed polynomials. With
$q=d+b+k+2\ge31$, one has $h\le q^6$, $h+1\le q^7$,
$Q_0\le q^6$, $r_{\max}\le q^7$, $H\le q^{64}$, and hence
\[
 R_1\le q^{71},\quad R_{\rm D}\le q^{82},\quad
 R_{\rm T}\le q^{285},\quad R\le q^{290}\le q^{300}.
\]
Here we used $T\le q^{202}$, $J_{\rm o}\le q^5$, and
$q_{\rm s}\le q^2$. Probability numerators never exceed their common
denominator. Intermediate products and sums may be stored with the
prescribed stage's denominator, so their bit lengths, including the
offset counts in \eqref{eq:offset-count}, are polynomial as well.
Candidate tables that fail feasibility can have negative dependent
entries, but their entries are integer sums of $O(d)$ quantities of size
$O(C)$ and still have $O(b+\log d)$ bits.

For each state, dense matrix propagation costs at most
$H\mathsf V^2$ arithmetic operations. Bin--offset summation, even with
a scan of all completions to identify the output, costs at most
$\mathsf V^3$; the geometric composition costs at most
$J_{\rm D}\mathsf V$. Constructing \eqref{eq:implemented-kernel} uses at
most $\mathsf S^2\mathsf V$ terms. Small-state propagation uses
$T\mathsf S^2$, terminal completion and output lookup at most
$\mathsf S\mathsf V^2$, and the last geometric composition at most
$J_{\rm o}\mathsf V$. Including the cache calculation for all
$\mathsf S$ states, the total number of arithmetic operations is at most
$2^{10db}\operatorname{poly}(d,b,k)$; multiplying by the established
polynomial bit cost per operation gives the stated, looser $2^{500db}$
bound. The iteration counts $H,J_{\rm D},T,J_{\rm o}$ depend polynomially
on $k$; none of the enumerated spaces does.
\end{proof}

\subsection{Exact uniform sampling}

The final step uses the rare residual-mixture construction of G\"obel,
Liu, Manurangsi, and Pappik~\cite[Section~3.1, Algorithm~1 and
Theorem~6]{GLMP2024}. A sufficiently accurate approximate law is dominated
pointwise by a slight enlargement of the target law. Its remaining mass
can therefore be sampled exactly on a rare branch. Here we verify this
domination, implement the mixture with unbiased bits, and use
Proposition~\ref{prop:law-tabulation} to account for the exhaustive work.

\begin{theorem}[Exact sampling]\label{thm:exact-sampling}
There is an algorithm using unbiased random bits that terminates almost
surely, outputs an exactly uniform element of $\Omega$, and has expected
bit running time polynomial in $d,b$.
\end{theorem}

\begin{proof}
Set
\[
 D=d^6b^2,\qquad k=2D,\qquad\delta=2^{-D},\qquad M=|\Omega|.
\]
Theorem~\ref{thm:almost-uniform} gives, for every $z\in\Omega$,
$p_k(z)\le1/M+\delta^2$. Also $M\le(C+1)^d\le2^{db}$ and $D\ge db$,
so $M\delta\le1$. Therefore
\[
 (1-\delta)p_k(z)
 \le\frac{1-\delta}{M}+(1-\delta)\delta^2
 \le\frac1M.
\]
It follows that
$r(z)=[1/M-(1-\delta)p_k(z)]/\delta$ is a probability law.

Use $D$ independent bits to enter a correction branch precisely when
all are zero. On the complementary event run the bounded sampler with
accuracy $k$. On the correction branch, enumerate $\Omega$ to obtain
$M$ and use Proposition~\ref{prop:law-tabulation} to compute
$p_k(z)=a_z/2^R$. Form the integer weights
\begin{equation}\label{eq:residual-integers}
 n_z=2^{R+D}-M(2^D-1)a_z,
 \qquad \sum_{z\in\Omega}n_z=M2^R.
\end{equation}
They are nonnegative by the pointwise domination just proved, and the
sum follows from $\sum_z a_z=2^R$. Thus $r(z)=n_z/(M2^R)$.
All integers in \eqref{eq:residual-integers} have
$O(R+D+\log M)$ bits; the fact that $M$ can be large as a number does not
make its binary representation long.

To sample this law exactly, put $V=M2^R$ and draw a uniform integer in
$\{0,\ldots,V-1\}$ by using $\lceil\log_2V\rceil$ bits and repeating
whenever their value is at least $V$. If $V=1$, use zero bits.
The acceptance probability is greater than $1/2$, so the expected number
of attempts is less than two, and termination is almost sure. A scan of
cumulative weights $n_z$ then selects the output exactly according to
$r$. Use fresh bits on each attempt and after the branch decision.

The output law is
$(1-\delta)p_k+\delta r=\operatorname{Uniform}(\Omega)$.
The correction branch has conditional expected cost at most
$2^{500db}\operatorname{poly}(d,b,k)$, including enumeration, integer
arithmetic, and the exact integer draw. Its contribution to the
unconditional expectation is polynomial, since
\[
 \delta\,2^{500db}
 =2^{-db(d^5b-500)}\le1
 \qquad(d\ge14,\ b\ge d).
\]
Substituting $k=2d^6b^2$ into a fixed polynomial preserves polynomial
dependence on $d,b$. The other branch and the branch decision also have
polynomial cost. This proves the expectation and termination claims.
\end{proof}

\begin{proof}[Proof of Theorem~\ref{thm:main}]
If either dimension is empty, if $N=0$, or if one dimension is one,
return the unique feasible table as in Section~\ref{sec:graph}.
Otherwise the parameters in \eqref{eq:parameters} satisfy
$d=O(mn+m+n+1)$ and $b=O(d+\log(N+1)+\log d)$.
For the supplied integer $k\ge1$, Theorem~\ref{thm:almost-uniform}
gives the first assertion. Theorem~\ref{thm:exact-sampling} gives the
second. Their bounds therefore have the stated dependence on dimensions
and binary margin size. All outputs are feasible, and the exact
algorithm's only unbounded loop is the almost-surely terminating integer
rejection on its correction branch.
\end{proof}
